\subsection{Formulas on the Calculator} 
When we plug unpleasant numbers into a formula on the calculator, it's very easy to mis-click. We can take advantage of the calculator's ability to store numbers to avoid re-typing values. Every graphing calculator will be able to do this, and some scientific calculators have the ability to store multiple values like this, too.
\begin{procedure}{Using Stored Values to Apply the Quadratic Formula}
\begin{enumerate*}
\item First, we will store our values of $a$, $b$, and $c$.
\begin{enumerate*}
\item Enter the value of $a$, press \TIbut{sto}, then press \TIbut{alpha}\ \TIbut{math} to enter the letter \TI{A}.
\item Repeat for $b$ (pressing \TIbut{alpha}\ \TIbut{apps} for the letter \TI{B})
\item Repeat for $c$ (pressing \TIbut{alpha}\ \TIbut{prgm} for the letter \TI{C})
\end{enumerate*}
\item Next, we'll compute the discriminant. Enter \TI{B\TIsup{2}-4AC}, then press \TIbut{sto}\ \TIbut{alpha}\ \TIbutpro{$x^{-1}$} to store this value as \TI{D}.
\item If $D\geq 0$, enter \TI{(-B+\textsurd D)/(2A)} to find the first solution, and \TI{(-B-\textsurd D)/(2A)} to find the second. Round appropriately.
\end{enumerate*}
\end{procedure}
\begin{example}
A company determines that its profit, in millions of dollars, from selling $x$ million widgets is given by the function
$$P(x)=-0.06x^2+3.8x$$
Use the quadratic formula to determine how many widgets they should sell to make a profit of \$50.1 million. Round your answers to one decimal place.
\begin{lpic}{}{7}
\regtextleft{0.25}{-6}{Check:};
\end{lpic}
\end{example}